% ============================================================
\setpartpreamble{\partintro
  {Metadata, Fluent API, conventions, validation and new column types.}
  {
    \item Chapter~\ref{ch:annotations-and-fluent-api}: Annotations and Fluent API
    \item Chapter~\ref{ch:conventions}: Conventions
    \item Chapter~\ref{ch:model-validation}: Model validation
    \item Chapter~\ref{ch:type-mappings-and-value-converters}: Type mappings and value
    converters
  }}
\part{The model}\label{part:model}

% ------------------------------------------------------------
\chapter{Annotations and Fluent API}\label{ch:annotations-and-fluent-api}
\versionrisk{stable}

Annotations are EF's mechanism for extensible metadata: key-value pairs on every model
element (model, entity type, property, index, ...). A plugin stores its configuration as
annotations, and the entire remaining infrastructure (conventions, migrations, snapshots,
compiled models) can work with them.

\section{Rules for annotations}

\begin{datatable}{H{0.75}H{1.25}}
  \textbf{Rule} & \textbf{Reason} \\
  \tablehead
  Own prefix (\code{Acme:Rls:}) & Avoids collisions with EF, providers and other plugins. \\
  \tablerule
  Simple values only: string, bool, numbers, enums & Snapshots and migrations write values
  as C\# literals. Custom classes cannot be serialized. \\
  \tablerule
  Complex values as a string (e.g. JSON) & The differ compares strings reliably via
  \code{Equals}. A record with a \code{string[]} would be compared by reference and
  constantly appear “changed”. \\
  \tablerule
  Do not store default values & “Annotation missing” and “annotation has default” differ for
  the differ and produce empty migrations. \\
  \tablerule
  One annotation per piece of configuration & Every policy has its own annotation, so the
  diff shows exactly which one changed. \\
\end{datatable}

\begin{sourcecode}[Metadata/\allowbreak{}Rls\allowbreak{}Annotation\allowbreak{}Names.\allowbreak{}cs,\allowbreak{} Metadata/\allowbreak{}Rls\allowbreak{}Policy.\allowbreak{}cs (\allowbreak{}shortened, simplified)]{csharp}
public static class RlsAnnotationNames
{
    public const string Prefix = "Acme:Rls:";
    public const string Enabled = Prefix + "Enabled";
    public const string Forced = Prefix + "Forced";
    public const string PolicyPrefix = Prefix + "Policy:";     // + policy name
    public const string QueryFilterKey = "Acme.Tenancy";
}

public sealed record RlsPolicy(string Name, PolicyCommand Command, bool Permissive,
    string[]? Roles, string? Using, string? WithCheck)
{
    public string Serialize()
        => JsonSerializer.Serialize(this, RlsJsonContext.Default.RlsPolicy);

    public static RlsPolicy Deserialize(string json) => /* ... */;
}
\end{sourcecode}

\section{The Fluent API}

Fluent API methods are extension methods on EF's builders. Three design rules make them
professional: the \textbf{non-generic} variant is the actual implementation (the model
snapshot uses non-generic builders, chapter~\ref{ch:design-time-services}). The
\textbf{generic} variant only delegates and preserves the type for further chaining. And the
class lives in the namespace
\code{Microsoft.\allowbreak{}Entity\allowbreak{}Framework\allowbreak{}Core} so it is
available everywhere without a \code{using}.

% source: src/Acme.EntityFrameworkCore.Tenancy.PostgreSQL/Builders/RlsEntityTypeBuilderExtensions.cs
\begin{sourcecode}[Builders/\allowbreak{}Rls\allowbreak{}Entity\allowbreak{}Type\allowbreak{}Builder\allowbreak{}Extensions.\allowbreak{}cs (\allowbreak{}shortened)]{csharp}
namespace Microsoft.EntityFrameworkCore;

public static class RlsEntityTypeBuilderExtensions
{
    public static EntityTypeBuilder HasRowLevelSecurity(
        this EntityTypeBuilder builder, bool force = false)
    {
        ArgumentNullException.ThrowIfNull(builder);
        builder.Metadata.SetAnnotation(RlsAnnotationNames.Enabled, true);

        if (force)
            builder.Metadata.SetAnnotation(RlsAnnotationNames.Forced, true);
        else
            builder.Metadata.RemoveAnnotation(RlsAnnotationNames.Forced);   // no default!

        return builder;
    }

    public static EntityTypeBuilder HasPolicy(
        this EntityTypeBuilder builder, string name, Action<PolicyBuilder> configure)
    {
        var policyBuilder = new PolicyBuilder(name);
        configure(policyBuilder);
        builder.Metadata.SetAnnotation(
            RlsAnnotationNames.PolicyPrefix + name, policyBuilder.Build().Serialize());
        return builder;
    }

    public static EntityTypeBuilder<TEntity> HasRowLevelSecurity<TEntity>(
        this EntityTypeBuilder<TEntity> builder, bool force = false) where TEntity : class
        => (EntityTypeBuilder<TEntity>)HasRowLevelSecurity(
            (EntityTypeBuilder)builder, force);

    // ... HasNoRowLevelSecurity, HasPolicy<TEntity> analogous
}
\end{sourcecode}

Nested configuration such as \code{Has\allowbreak{}Policy} accepts an
\code{Action<\allowbreak{}Policy\allowbreak{}Builder>}. The builder checks the database's
rules already while the model is built. A mistake like “INSERT policy with USING” therefore
surfaces on the application's first start, not when the migration runs in production:

% source: src/Acme.EntityFrameworkCore.Tenancy.PostgreSQL/Builders/PolicyBuilder.cs (strings)
\begin{sourcecode}[Builders/\allowbreak{}Policy\allowbreak{}Builder.\allowbreak{}cs › Build]{csharp}
internal RlsPolicy Build()
{
    if (_using is null && _withCheck is null)
        throw new InvalidOperationException(
            $"Policy '{_name}' requires USING or WITH CHECK.");

    if (_command == PolicyCommand.Insert && _using is not null)
        throw new InvalidOperationException(
            $"INSERT policy '{_name}' only allows WITH CHECK.");

    if (_command is PolicyCommand.Select or PolicyCommand.Delete && _withCheck is not null)
        throw new InvalidOperationException(
            $"{_command} policy '{_name}' only allows USING.");

    return new RlsPolicy(_name, _command, _permissive,
        _roles.Count > 0 ? [.. _roles] : null, _using, _withCheck);
}
\end{sourcecode}

\section{Configuration sources}

EF remembers for every piece of metadata \emph{where} it came from. A source can override an
equally strong or weaker one, but never a stronger one. This is the mechanism through which
plugins set sensible defaults that users can always override:

\begin{datatable}{H{0.82}H{1.64}H{0.55}}
  \textbf{ConfigurationSource} & \textbf{Set by} & \textbf{Strength} \\
  \tablehead
  \code{Explicit} & Fluent API in \code{On\allowbreak{}Model\allowbreak{}Creating}
  (\code{Metadata.\allowbreak{}Set\allowbreak{}Annotation}) & strongest \\
  \tablerule
  \code{Data\allowbreak{}Annotation} & Attribute conventions
  (\code{from\allowbreak{}Data\allowbreak{}Annotation: true}) & medium \\
  \tablerule
  \code{Convention} & Conventions via \code{IConvention*Builder} & weakest \\
\end{datatable}

For custom conventions the plugin therefore also offers overloads on the convention builder.
They return \code{null} when a stronger source decided otherwise:

% source: src/Acme.EntityFrameworkCore.Tenancy.PostgreSQL/Builders/RlsEntityTypeBuilderExtensions.cs
\begin{sourcecode}[Builders/\allowbreak{}Rls\allowbreak{}Entity\allowbreak{}Type\allowbreak{}Builder\allowbreak{}Extensions.\allowbreak{}cs › convention variant]{csharp}
public static IConventionEntityTypeBuilder? HasRowLevelSecurity(
    this IConventionEntityTypeBuilder builder,
    bool enabled,
    bool fromDataAnnotation = false)
    => builder.HasAnnotation(RlsAnnotationNames.Enabled, enabled, fromDataAnnotation)
           is null
        ? null
        : builder;
\end{sourcecode}

The check of the source lives in the builder, not in the metadata. A convention that writes
directly to the metadata bypasses it:

\begin{sourcecode}[Convention (\allowbreak{}WRONG)]{csharp}
// In ProcessModelFinalizing: directly on the metadata instead of through entityType.Builder
entityType.SetAnnotation(RlsAnnotationNames.Enabled, true);
\end{sourcecode}

\begin{failure}[What happens if a convention bypasses the configuration source?]
  \textbf{Observable.} The user switches RLS off for an entity with
  \code{Has\allowbreak{}No\allowbreak{}Row\allowbreak{}Level\allowbreak{}Security(\allowbreak{})},
  and after model building it is on again. The next migration enables RLS for a table that
  was explicitly excluded.

  \textbf{Cause.} \code{Set\allowbreak{}Annotation} on the convention metadata overwrites
  any existing value. Only the convention builder compares the sources and returns
  \code{null} when a stronger one has decided.

  \textbf{Solution.} In conventions always write through
  \code{entity\allowbreak{}Type.\allowbreak{}Builder} or use the plugin's convention
  overloads.

  \textbf{Test:}
  \code{Failure\allowbreak{}Case\allowbreak{}Tests.\allowbreak{}Convention\allowbreak{}\_writing\allowbreak{}\_to\allowbreak{}\_the\allowbreak{}\_metadata\allowbreak{}\_…}
  and the counter-check through the builder.
\end{failure}

% ------------------------------------------------------------
\chapter{Conventions}\label{ch:conventions}
\versionrisk{stable}

Conventions change the model automatically. EF itself consists largely of conventions:
discovering keys, finding relationships, deriving table names. A plugin adds its own to make
configuration unnecessary.

\section{Convention types and timing}

A convention implements one or more interfaces. Each stands for an event during model
building. The most important ones:

\begin{datatable}{H{1.04}H{0.91}H{1.04}}
  \textbf{Interface} & \textbf{Event} & \textbf{Typical use} \\
  \tablehead
  \code{IEntity\allowbreak{}Type\allowbreak{}Added\allowbreak{}Convention} & Entity type was
  added & Evaluate attributes, base configuration \\
  \tablerule
  \code{IProperty\allowbreak{}Added\allowbreak{}Convention} & Property was added & Column
  types, naming conventions \\
  \tablerule
  \code{IForeign\allowbreak{}Key\allowbreak{}Added\allowbreak{}Convention} & Relationship
  created & Indexes, delete behavior \\
  \tablerule
  \code{IEntity\allowbreak{}Type\allowbreak{}Base\allowbreak{}Type\allowbreak{}Changed\allowbreak{}Convention}
  & Inheritance changed & Move configuration \\
  \tablerule
  \code{IModel\allowbreak{}Finalizing\allowbreak{}Convention} & Model is being finalized
  (still mutable) & Everything that depends on the finished model \\
  \tablerule
  \code{IModel\allowbreak{}Finalized\allowbreak{}Convention} & Model is finished (read-only)
  & Validation \\
\end{datatable}

The most common convention mistake is the wrong point in time:

\begin{sourcecode}[Convention (\allowbreak{}WRONG)]{csharp}
public void ProcessPropertyAdded(IConventionPropertyBuilder propertyBuilder,
    IConventionContext<IConventionPropertyBuilder> context)
{
    // Remember the column name right away, for example for the policy expression
    propertyBuilder.HasAnnotation("Acme:Column", propertyBuilder.Metadata.GetColumnName());
}
\end{sourcecode}

\begin{failure}[What happens if a convention runs too early?]
  \textbf{Observable.} If the user then configures
  \code{Has\allowbreak{}Column\allowbreak{}Name(\textquotedbl{}tenant\textquotedbl{})} or a
  naming plugin renames the column, the remembered name stays
  \code{Tenant\allowbreak{}Id}. The policy refers to a column that does not exist, and the
  migration only fails in the database.

  \textbf{Cause.} Reactive conventions (\code{*Added}, \code{*Changed}) run while the user is
  still configuring. A column name may still change afterwards, an entity type may still get
  a base type.

  \textbf{Solution.} Read whatever depends on the finished model in an
  \code{IModel\allowbreak{}Finalizing\allowbreak{}Convention}, through
  \code{Get\allowbreak{}Column\allowbreak{}Name(\allowbreak{}table)}. A convention that must
  stay reactive reacts to \emph{all} relevant events.

  \textbf{Test:}
  \code{Failure\allowbreak{}Case\allowbreak{}Tests.\allowbreak{}Convention\allowbreak{}\_reading\allowbreak{}\_the\allowbreak{}\_column\allowbreak{}\_name\allowbreak{}\_too\allowbreak{}\_early\allowbreak{}\_…}
  and the counter-check at \code{Model\allowbreak{}Finalizing}.
\end{failure}

\section{Plugging in conventions}

The plugin interface \code{IConvention\allowbreak{}Set\allowbreak{}Plugin} receives the
convention set and adds its own conventions.
\code{Convention\allowbreak{}Set.\allowbreak{}Add} registers a convention in all lists whose
interfaces it implements, each time at the end:

% source: src/Acme.EntityFrameworkCore.Tenancy.PostgreSQL/Conventions/TenancyConventionSetPlugin.cs
\begin{sourcecode}[Conventions/\allowbreak{}Tenancy\allowbreak{}Convention\allowbreak{}Set\allowbreak{}Plugin.\allowbreak{}cs]{csharp}
public class TenancyConventionSetPlugin : IConventionSetPlugin
{
    private readonly ProviderConventionSetBuilderDependencies _dependencies;
    private readonly IDbContextOptions _options;
    private readonly ISqlGenerationHelper _sqlGenerationHelper;
    private readonly TenancyLoggingDefinitions _loggingDefinitions;

    public TenancyConventionSetPlugin(
        ProviderConventionSetBuilderDependencies dependencies,
        IDbContextOptions options,
        ISqlGenerationHelper sqlGenerationHelper,
        TenancyLoggingDefinitions loggingDefinitions)
    {
        // ... ArgumentNullException.ThrowIfNull for all parameters ...
        _dependencies = dependencies;
        _options = options;
        _sqlGenerationHelper = sqlGenerationHelper;
        _loggingDefinitions = loggingDefinitions;
    }

    public virtual ConventionSet ModifyConventions(ConventionSet conventionSet)
    {
        var settings = _options.FindExtension<TenancyOptionsExtension>()?.Settings;
        if (settings is null)
            return conventionSet;

        // chapter "Value generation"
        conventionSet.Add(new TenantIdKeyValueGenerationConvention());
        conventionSet.Add(new NoTenantIsolationAttributeConvention(_dependencies));
        conventionSet.Add(new TenantIsolationConvention(
            settings, _sqlGenerationHelper, _dependencies.Logger, _loggingDefinitions));
        conventionSet.Add(new TenancyModelValidationConvention(
            _dependencies.Logger, _loggingDefinitions));

        return conventionSet;
    }
}
\end{sourcecode}

\code{Provider\allowbreak{}Convention\allowbreak{}Set\allowbreak{}Builder\allowbreak{}Dependencies}
is the parameter object that EF's own conventions receive too. Among other things it
contains the model logger, the validation logger and the type mapping source. Other services
such as \code{ISql\allowbreak{}Generation\allowbreak{}Helper} can be injected as well,
because the plugin itself is resolved from the internal provider.

\section{The isolation convention}

The sample's central convention runs at \code{Model\allowbreak{}Finalizing} because column
names are only final then. It finds the tenant property (marked with
\code{Is\allowbreak{}Tenant\allowbreak{}Property(\allowbreak{})} or
\code{Tenant\allowbreak{}Id} for \code{ITenant\allowbreak{}Owned},
chapter~\ref{ch:property-level-fluent-api-and}), enables RLS with a policy, creates an
index, sets a runtime annotation for the SaveChanges interceptor and optionally a query
filter:

\begin{sourcecode}[Conventions/\allowbreak{}Tenant\allowbreak{}Isolation\allowbreak{}Convention.\allowbreak{}cs › Process\allowbreak{}Model\allowbreak{}Finalizing (\allowbreak{}simplified)]{csharp}
public virtual void ProcessModelFinalizing(
    IConventionModelBuilder modelBuilder,
    IConventionContext<IConventionModelBuilder> context)
{
    // TPT, TPC, entity splitting, owned types in their own table and implicit join
    // tables would leave tenant data without a tenant column: reject them first
    TenantTableMappingValidator.Validate(modelBuilder.Metadata, e => e.IsRlsOptedOut());

    foreach (var entityType in modelBuilder.Metadata.GetEntityTypes())
    {
        if (entityType.BaseType is not null       // TPH: the root type's table
            || entityType.IsOwned())              // validated: owner's table or JSON
            continue;

        // Marked property or "TenantId" with ITenantOwned (core package,
        // chapter "Property-level Fluent API and runtime annotations")
        if (entityType.FindTenantPropertyCandidate() is not IConventionProperty property)
            continue;                             // (+ type check, shortened)

        if (entityType.IsRlsOptedOut())
        {
            _logger.TenantIsolationOptedOut(_loggingDefinitions, entityType);  // warning
            continue;
        }

        // No table (view, ToSqlQuery, keyless) -> no RLS
        if (StoreObjectIdentifier.Create(entityType, StoreObjectType.Table)
            is not { } table)
            continue;

        // Runtime annotation: this is how the SaveChanges interceptor finds the property
        entityType.Builder.HasAnnotation(
            TenancyAnnotationNames.TenantProperty, property.Name);

        ApplyRowLevelSecurity(entityType, property, table);
        if (_settings.CreateTenantIndex) EnsureTenantIndex(entityType, property);
        if (_settings.ApplyQueryFilter) ApplyQueryFilter(entityType, property);

        _logger.TenantIsolationApplied(_loggingDefinitions, entityType);       // debug
    }
}
\end{sourcecode}

% source: src/Acme.EntityFrameworkCore.Tenancy.PostgreSQL/Conventions/TenantIsolationConvention.cs
\begin{sourcecode}[Conventions/\allowbreak{}Tenant\allowbreak{}Isolation\allowbreak{}Convention.\allowbreak{}cs › Apply\allowbreak{}Row\allowbreak{}Level\allowbreak{}Security]{csharp}
protected virtual void ApplyRowLevelSecurity(
    IConventionEntityType entityType,
    IConventionProperty property,
    StoreObjectIdentifier table)
{
    // Column name for exactly this table (may differ with table splitting)
    var column = _sqlGenerationHelper.DelimitIdentifier(property.GetColumnName(table)!);
    var condition =
        $"{column} = nullif(current_setting('{_settings.SettingName}', true), '')::uuid";

    var policy = new RlsPolicy(_settings.PolicyName, PolicyCommand.All, Permissive: true,
        Roles: null, Using: condition, WithCheck: condition);

    // Convention builder: does NOT override explicit configuration
    var builder = entityType.Builder;
    builder.HasAnnotation(RlsAnnotationNames.Enabled, true);
    if (_settings.ForceRowLevelSecurity)
        builder.HasAnnotation(RlsAnnotationNames.Forced, true);
    builder.HasAnnotation(RlsAnnotationNames.PolicyPrefix + _settings.PolicyName,
        policy.Serialize());
}
\end{sourcecode}

\begin{important}[Validate SQL you generate yourself]
  The setting name is written into the policy expression without escaping. Unlike SQL that
  the developer writes in
  \code{Using(\allowbreak{}.\allowbreak{}.\allowbreak{}.\allowbreak{})} themselves, here the
  \emph{plugin} assembles SQL and is responsible for keeping invalid input out. That is why
  the options builder validates the name with a strict regular expression
  (\code{[a-z\_\allowbreak{}][a-z0-9\_\allowbreak{}]*\textbackslash{}.\allowbreak{}[a-z\_\allowbreak{}][a-z0-9\_\allowbreak{}]*}).
\end{important}

\section{Which table mappings are protected}\label{sec:supported-table-mappings}

RLS protects tables, not entities. The policy compares the tenant column of the \emph{same} table
with the session setting. A table that holds tenant data but lacks this column cannot be
protected by this policy. One entity can occupy several tables, and several entities can share
one table. That is why \code{Tenant\allowbreak{}Table\allowbreak{}Mapping\allowbreak{}Validator}
from the relational package checks, before any annotation, that every table with tenant data contains
the tenant column. Both provider packages call it, so the contract is the same for PostgreSQL and
SQL Server:

\begin{datatable}{H{1.00}H{1.25}H{0.75}}
  \textbf{Mapping} & \textbf{Tables with tenant data} & \textbf{Behavior} \\
  \tablehead
  One table per entity & the table, with tenant column & protected \\
  \tablerule
  TPH (default) & one table for the whole hierarchy & protected, policy on the root type \\
  \tablerule
  Tenant-owned derived type below a root without a tenant & the same table, rows of other types
  without a tenant & rejected \\
  \tablerule
  TPT & one table per derived type \emph{without} tenant column & rejected \\
  \tablerule
  TPC & one table per concrete type, the abstract base has none & rejected (not
  implemented) \\
  \tablerule
  Owned type in the owner's table (\code{Owns\allowbreak{}One}) & columns of the owner's
  table & protected \\
  \tablerule
  Owned type as JSON (\code{To\allowbreak{}Json}) & one column of the owner's table &
  protected \\
  \tablerule
  Owned type in its own table (\code{Owns\allowbreak{}Many}, \code{To\allowbreak{}Table}) &
  own table without tenant column & rejected \\
  \tablerule
  Table splitting & the same rows for all types & protected if all tenant-owned types use the
  same column \\
  \tablerule
  Entity splitting (\code{Split\allowbreak{}To\allowbreak{}Table}) & another table without
  tenant column & rejected \\
  \tablerule
  Implicit many-to-many table with a tenant-owned side & join table without tenant column &
  rejected \\
  \tablerule
  Join entity with a tenant property (\code{Using\allowbreak{}Entity}) & join table with
  tenant column & protected \\
  \tablerule
  View, \code{To\allowbreak{}Sql\allowbreak{}Query}, keyless without a table & no table of
  its own & no policy; the base tables carry it \\
\end{datatable}

Rejected means: the model throws an \code{Invalid\allowbreak{}Operation\allowbreak{}Exception}
on the first access to \code{Model}, naming the table, the type and the way out. This happens
before \code{Ensure\allowbreak{}Created}, before any migration and before the first query.
Seeming support that writes a policy onto a table without a matching column would be worse than
the rejection. TPC could be protected in principle, since every concrete table inherits the
tenant column. The plugin, however, supports exactly one table per hierarchy, and a half-done
implementation must not pretend to be secure. The ways out are TPH or independent entities.

Deliberate exceptions remain possible: an entity excluded with
\code{[No\allowbreak{}Tenant\allowbreak{}Isolation]} or
\code{Has\allowbreak{}No\allowbreak{}Row\allowbreak{}Level\allowbreak{}Security(\allowbreak{})}
is not checked, and the convention logs the exception as a warning. That way the model tells a
decision from an oversight. For the implicit join table the PostgreSQL package offers
\code{Using\allowbreak{}Entity(\allowbreak{}j => j.\allowbreak{}Has\allowbreak{}No\allowbreak{}Row\allowbreak{}Level\allowbreak{}Security(\allowbreak{}))};
the SQL Server package only knows the attribute and needs a join class for it. A tenant property
on the join entity is almost always better:

\begin{sourcecode}[Application]{csharp}
modelBuilder.Entity<Basket>()
    .HasMany(b => b.Labels).WithMany(l => l.Baskets)
    .UsingEntity("BasketLabel",
        j => j.Property<TenantId>("TenantId").IsTenantProperty());   // shadow property
\end{sourcecode}

The SaveChanges interceptor fills the shadow property like any other tenant property.
\code{Mapping\allowbreak{}Isolation\allowbreak{}Tests} checks both against a real PostgreSQL
with a restricted role: a TPC model fails before a single table exists. In a model with TPH, an
owned type, a JSON column and this join entity, every table has RLS, and without a tenant no row
is visible in any table. \code{Table\allowbreak{}Mapping\allowbreak{}Tests} and
\code{Sql\allowbreak{}Server\allowbreak{}Table\allowbreak{}Mapping\allowbreak{}Tests} cover
every row of the matrix without a database.

\section{Structural changes: index and query filter}

Conventions may not only set annotations but change the model structurally. The index is
only created if no index or key already starts with the tenant column. The query filter is
built as an expression tree and registered as a \textbf{named filter}, new in EF Core 10.
Named filters coexist with the user's own filters and can be switched off selectively with
\code{Ignore\allowbreak{}Query\allowbreak{}Filters(\allowbreak{}[\textquotedbl{}Acme.\allowbreak{}Tenancy\textquotedbl{}])}.

% source: src/Acme.EntityFrameworkCore.Tenancy.PostgreSQL/Conventions/TenantIsolationConvention.cs
\begin{sourcecode}[Conventions/\allowbreak{}Tenant\allowbreak{}Isolation\allowbreak{}Convention.\allowbreak{}cs (\allowbreak{}excerpt)]{csharp}
protected virtual void EnsureTenantIndex(
    IConventionEntityType entityType, IConventionProperty property)
{
    var covered = entityType.GetIndexes().Any(i => i.Properties[0] == property)
                  || entityType.GetKeys().Any(k => k.Properties[0] == property);
    if (!covered)
        entityType.Builder.HasIndex(new[] { property });
}

protected virtual void ApplyQueryFilter(
    IConventionEntityType entityType, IConventionProperty property)
{
    // e => EF.Property<TenantId>(e, "...") == EF.Functions.CurrentTenant()
    // EF.Property works for any property name and for shadow properties too
    var parameter = Expression.Parameter(entityType.ClrType, "e");
    var body = Expression.Equal(
        Expression.Call(typeof(EF), nameof(EF.Property), [typeof(TenantId)],
            parameter, Expression.Constant(property.Name)),
        Expression.Call(CurrentTenantMethod, Expression.Property(null, FunctionsProperty)));

    entityType.Builder.HasQueryFilter(
        RlsAnnotationNames.QueryFilterKey, Expression.Lambda(body, parameter));
}
\end{sourcecode}

Notably, the filter uses the plugin's own function
\code{EF.\allowbreak{}Functions.\allowbreak{}Current\allowbreak{}Tenant(\allowbreak{})}. It
is therefore translated to SQL by the translator from
chapter~\ref{ch:translating-methods-and-members}. The plugin's building blocks interlock.

\textbf{Combined with soft delete.} Before EF Core~10 an entity type had exactly one query
filter: a second \code{Has\allowbreak{}Query\allowbreak{}Filter} replaced the first, and
\code{Ignore\allowbreak{}Query\allowbreak{}Filters(\allowbreak{})} switched off all of them at
once, the tenant filter included. With named filters the application's own filters stay
independent of the plugin's:

% source: tests/Acme.EntityFrameworkCore.Tenancy.Tests/QueryTranslationTests.cs
\begin{sourcecode}[tests/\allowbreak{}.\allowbreak{}.\allowbreak{}.\allowbreak{}/\allowbreak{}Query\allowbreak{}Translation\allowbreak{}Tests.\allowbreak{}cs (\allowbreak{}excerpt)]{csharp}
m.Entity<Order>().HasQueryFilter("SoftDelete", o => !EF.Property<bool>(o, "IsDeleted"));

// WHERE NOT (o."IsDeleted") AND o."TenantId" = nullif(current_setting(...), '')::uuid
var both = ctx.Orders.ToQueryString();

// Deleted rows become visible, the tenant filter stays
var withDeleted = ctx.Orders.IgnoreQueryFilters(["SoftDelete"]).ToQueryString();

// Only the tenant filter is off (with RLS the database still filters)
var allTenants = ctx.Orders.IgnoreQueryFilters(["Acme.Tenancy"]).ToQueryString();
\end{sourcecode}

\begin{important}[Switch off filters by name only]
  An \code{Ignore\allowbreak{}Query\allowbreak{}Filters(\allowbreak{})} without names also
  switches off the tenant filter. Where an application shows deleted records, for example in a
  recycle bin, it names its own filter:
  \code{Ignore\allowbreak{}Query\allowbreak{}Filters(\allowbreak{}[\textquotedbl{}Soft\allowbreak{}Delete\textquotedbl{}])}.
  Switching off \code{Acme.\allowbreak{}Tenancy} belongs to support or reporting roles only:
  with RLS the database still filters, without it (Cosmos DB,
  chapter~\ref{ch:without-sql-what-carries-over}) the query then sees every tenant.
\end{important}

\section{Final names, whoever sets them}

The isolation convention reads the column name only at \code{Model\allowbreak{}Finalizing}
and only through \code{Get\allowbreak{}Column\allowbreak{}Name(\allowbreak{}table)}. Names
that other conventions or plugins change are final by then. The sample checks this against
EFCore.NamingConventions, which renames tables and columns to snake\_case, in both
activation orders:

\begin{sourcecode}[Usage]{csharp}
options.UseNpgsql(connectionString).UseAcmeTenancy().UseSnakeCaseNamingConvention();
\end{sourcecode}

\begin{sourcecode}[Generated migration SQL (\allowbreak{}excerpt,\allowbreak{} tests/\allowbreak{}.\allowbreak{}.\allowbreak{}.\allowbreak{}/\allowbreak{}Naming\allowbreak{}Convention\allowbreak{}Tests.\allowbreak{}cs)]{sql}
ALTER TABLE orders ENABLE ROW LEVEL SECURITY;
CREATE POLICY tenant_isolation ON orders AS PERMISSIVE FOR ALL
    USING (tenant_id = nullif(current_setting('app.tenant_id', true), '')::uuid)
    WITH CHECK (tenant_id = nullif(current_setting('app.tenant_id', true), '')::uuid);
CREATE INDEX ix_orders_tenant_id ON orders (tenant_id);
\end{sourcecode}

\begin{tip}[Never assemble names yourself]
  A plugin that assembles names itself (from \code{Name} properties or by string
  concatenation) breaks as soon as a naming plugin is active. Always ask the model for the
  final name, and as late as possible.
\end{tip}

\section{Attribute conventions}

Data annotations are evaluated with an attribute convention. EF provides the base class
\code{Type\allowbreak{}Attribute\allowbreak{}Convention\allowbreak{}Base<\allowbreak{}TAttribute>},
which finds the attribute and only asks for the reaction. The value is set with
\code{from\allowbreak{}Data\allowbreak{}Annotation: true} so the Fluent API can override it:

% source: src/Acme.EntityFrameworkCore.Tenancy.PostgreSQL/Conventions/NoTenantIsolationAttributeConvention.cs
\begin{sourcecode}[Conventions/\allowbreak{}No\allowbreak{}Tenant\allowbreak{}Isolation\allowbreak{}Attribute\allowbreak{}Convention.\allowbreak{}cs]{csharp}
public class NoTenantIsolationAttributeConvention
    : TypeAttributeConventionBase<NoTenantIsolationAttribute>
{
    public NoTenantIsolationAttributeConvention(
        ProviderConventionSetBuilderDependencies dependencies)
        : base(dependencies)
    {
    }

    protected override void ProcessEntityTypeAdded(
        IConventionEntityTypeBuilder entityTypeBuilder,
        NoTenantIsolationAttribute attribute,
        IConventionContext<IConventionEntityTypeBuilder> context)
        => entityTypeBuilder.HasRowLevelSecurity(enabled: false, fromDataAnnotation: true);
}
\end{sourcecode}

\begin{tip}
  The attribute itself lives in the \emph{core package}, the convention in the provider
  package. Entity classes can use the attribute without depending on PostgreSQL. This is the
  same split as with EF's own attributes in
  \code{System.\allowbreak{}Component\allowbreak{}Model.\allowbreak{}Data\allowbreak{}Annotations}.
\end{tip}

% ------------------------------------------------------------
\chapter{Model validation}\label{ch:model-validation}
\versionrisk{stable}

A plugin should detect misconfiguration before it leads to wrong SQL or data leaks. EF
validates models through \code{IModel\allowbreak{}Validator}, but there is only one, and it
belongs to the provider. For plugins, an
\code{IModel\allowbreak{}Finalized\allowbreak{}Convention} is the better place: it runs on
the finished model and collides with nobody.

\begin{sourcecode}[Conventions/\allowbreak{}Tenancy\allowbreak{}Model\allowbreak{}Validation\allowbreak{}Convention.\allowbreak{}cs (\allowbreak{}simplified)]{csharp}
public virtual IModel ProcessModelFinalized(IModel model)
{
    foreach (var entityType in model.GetEntityTypes())
    {
        if (!entityType.HasRlsConfiguration())
            continue;

        if (entityType.BaseType is not null)
            throw new InvalidOperationException(
                $"RLS on '{entityType.DisplayName()}' is invalid: only allowed on the root "
                + $"type ('{entityType.GetRootType().DisplayName()}').");

        if (StoreObjectIdentifier.Create(entityType, StoreObjectType.Table) is null)
            throw new InvalidOperationException("... not mapped to a table ...");

        if (!entityType.IsRlsEnabled() && entityType.GetRlsPolicies().Count > 0)
            _logger.PoliciesWithoutRowLevelSecurity(_loggingDefinitions, entityType);
    }
    return model;
}
\end{sourcecode}

The validation distinguishes two kinds of problems. \textbf{Structural errors} that are
guaranteed to lead to wrong behavior throw an exception. \textbf{Questionable configuration}
that may be intended produces a warning through the logging system. Users can turn such
warnings into errors selectively (chapter~\ref{ch:diagnostics-and-logging}).

% ------------------------------------------------------------
\chapter{Type mappings and value converters}\label{ch:type-mappings-and-value-converters}
\versionrisk{medium}

A type mapping connects a CLR type with a database type. It determines the column type in
migrations, how parameters are set, what literals look like in SQL, how values are compared
and copied and how they are stored in JSON columns. For custom types a plugin has two
options:

\begin{datatable}{H{0.92}H{1.04}H{1.04}}
  \textbf{Approach} & \textbf{Advantages} & \textbf{Disadvantages} \\
  \tablehead
  Value converter (user configures \code{Has\allowbreak{}Conversion}) & No plugin needed &
  Must be configured for every property or via \code{Configure\allowbreak{}Conventions} \\
  \tablerule
  \code{IRelational\allowbreak{}Type\allowbreak{}Mapping\allowbreak{}Source\allowbreak{}Plugin}
  & Type works everywhere without configuration, own literals, full control &
  Provider-specific (store type) \\
\end{datatable}

\section{The type and its converter}

The strongly typed \code{Tenant\allowbreak{}Id} and its converter live in the core package.
The converter has a shared instance; the section on compiled models explains why.

\begin{sourcecode}[Tenant\allowbreak{}Id.\allowbreak{}cs,\allowbreak{} Tenant\allowbreak{}Id\allowbreak{}Converter.\allowbreak{}cs (\allowbreak{}shortened, simplified)]{csharp}
public readonly record struct TenantId(Guid Value)
{
    public static TenantId Empty => default;
    public bool IsEmpty => Value == Guid.Empty;              // translated to SQL
    public static TenantId New() => new(Guid.CreateVersion7());
}

public sealed class TenantIdConverter : ValueConverter<TenantId, Guid>
{
    public static TenantIdConverter Instance { get; } = new();

    public TenantIdConverter()
        : base(id => id.Value, value => new TenantId(value))
    {
    }
}
\end{sourcecode}

\section{The type mapping}

\begin{sourcecode}[Storage/\allowbreak{}Tenant\allowbreak{}Id\allowbreak{}Type\allowbreak{}Mapping.\allowbreak{}cs (\allowbreak{}simplified)]{csharp}
public class TenantIdTypeMapping : RelationalTypeMapping
{
    // Required for compiled models: public static Default property
    public static TenantIdTypeMapping Default { get; } = new();

    public TenantIdTypeMapping()
        : base(new RelationalTypeMappingParameters(
            new CoreTypeMappingParameters(
                typeof(TenantId),
                converter: TenantIdConverter.Instance,
                jsonValueReaderWriter: TenantIdJsonReaderWriter.Instance),
            storeType: "uuid",
            dbType: System.Data.DbType.Guid))
    {
    }

    protected TenantIdTypeMapping(RelationalTypeMappingParameters parameters)
        : base(parameters)
    {
    }

    // EF clones mappings, e.g. for facets or changed comparers
    protected override RelationalTypeMapping Clone(
        RelationalTypeMappingParameters parameters)
        => new TenantIdTypeMapping(parameters);

    protected override string GenerateNonNullSqlLiteral(object value)
    {
        var guid = value switch
        {
            TenantId id => id.Value,
            Guid g => g,
            _ => throw new ArgumentException($"Unexpected type {value.GetType()}."),
        };
        return $"'{guid}'::uuid";
    }
}
\end{sourcecode}

Every custom mapping needs the protected constructor taking
\code{Relational\allowbreak{}Type\allowbreak{}Mapping\allowbreak{}Parameters} and a
\code{Clone} override that creates its own type. Without it, a cloned mapping loses its
behavior, for example literal generation.

\section{The type mapping plugin}

\code{Relational\allowbreak{}Type\allowbreak{}Mapping\allowbreak{}Source} asks all plugins
\emph{before} falling back to the provider's mappings. A plugin returns \code{null} when it
is not responsible:

\begin{sourcecode}[Storage/\allowbreak{}Tenancy\allowbreak{}Type\allowbreak{}Mapping\allowbreak{}Source\allowbreak{}Plugin.\allowbreak{}cs (\allowbreak{}simplified)]{csharp}
public class TenancyTypeMappingSourcePlugin : IRelationalTypeMappingSourcePlugin
{
    public virtual RelationalTypeMapping? FindMapping(
        in RelationalTypeMappingInfo mappingInfo)
    {
        if (mappingInfo.ClrType != typeof(TenantId))
            return null;

        // Explicitly different column type (HasColumnType("text")): not responsible
        var storeTypeBase = mappingInfo.StoreTypeNameBase;
        if (storeTypeBase is not null
            && !string.Equals(storeTypeBase, "uuid", StringComparison.OrdinalIgnoreCase))
            return null;

        return TenantIdTypeMapping.Default;
    }
}
\end{sourcecode}

\begin{note}
  Without a mapping, EF would not recognize \code{Tenant\allowbreak{}Id} as a scalar type
  during model building and would abort with an error. With the plugin, the type works as a
  property, as a key, in queries, as a parameter and in JSON columns, without a single line
  of configuration.
\end{note}

\section{JSON support}

For JSON columns a mapping needs a
\code{Json\allowbreak{}Value\allowbreak{}Reader\allowbreak{}Writer}. It has three duties:
read, write, and describe how it is created in generated code. The pattern follows EF's own
\code{Json\allowbreak{}Guid\allowbreak{}Reader\allowbreak{}Writer}:

% source: src/Acme.EntityFrameworkCore.Tenancy/TenantIdJsonReaderWriter.cs
\begin{sourcecode}[Tenant\allowbreak{}Id\allowbreak{}Json\allowbreak{}Reader\allowbreak{}Writer.\allowbreak{}cs]{csharp}
public sealed class TenantIdJsonReaderWriter : JsonValueReaderWriter<TenantId>
{
    private static readonly PropertyInfo InstanceProperty =
        typeof(TenantIdJsonReaderWriter).GetProperty(nameof(Instance))!;

    public static TenantIdJsonReaderWriter Instance { get; } = new();

    private TenantIdJsonReaderWriter()
    {
    }

    public override TenantId FromJsonTyped(
        ref Utf8JsonReaderManager manager, object? existingObject = null)
        => new(manager.CurrentReader.GetGuid());

    public override void ToJsonTyped(Utf8JsonWriter writer, TenantId value)
        => writer.WriteStringValue(value.Value);

    // For compiled models and precompiled queries
    public override Expression ConstructorExpression
        => Expression.Property(null, InstanceProperty);
}
\end{sourcecode}

\begin{important}[Compiled models need Default and Instance]
  When generating for NativeAOT (\code{dotnet ef dbcontext optimize -{}-nativeaot}), the
  compiled model generator (chapter~\ref{ch:compiled-models-trimming-and-nativeaot}) emits
  code of the form
  \code{Tenant\allowbreak{}Id\allowbreak{}Type\allowbreak{}Mapping.\allowbreak{}Default.\allowbreak{}Clone(\allowbreak{}.\allowbreak{}.\allowbreak{}.\allowbreak{})}
  for every type mapping. It requires a \textbf{public static property} named \code{Default}
  on a \textbf{public} type and aborts otherwise. It only takes over converters and
  readers/writers without extra code when they are \emph{the same instance} as in the
  default mapping. That is why the mapping uses the static instances rather than \code{new
  .\allowbreak{}.\allowbreak{}.\allowbreak{}(\allowbreak{})}. Without NativeAOT, EF 10
  resolves type mappings at startup instead; a plugin should still meet the requirements,
  because users decide on NativeAOT, not the plugin.
\end{important}
